\section{Discussion}

\subsection{Physical interpretation}

A brief early Si-rich channel can explain the lower high-Si regime without immediately breaking the relic fit. The accepted model suggests enrichment confined to the first few tens of Myr, before the model accumulates too much total metal mass or suppresses \mgfe\ beyond the accepted limit.

\subsection{Si/Mg, Mg/Fe, and relicness}

Relicness should correlate robustly with \mgfe: old compact relics form rapidly and avoid late Fe-rich dilution. Si/Mg is a more conditional relicness tracer. It may correlate with relicness only if an early Si-rich channel operated and if later accreted or extended star formation did not dilute the signal.

\subsection{Why normal ETGs may dilute toward ordinary Si/Mg}

Normal early-type galaxies can build more extended stellar populations through longer star-formation histories, mergers, and accretion. Those later populations can dilute an early high-Si component in both mass and light, pushing integrated \simg\ toward more ordinary values even if the earliest compact core was Si-rich.

\subsection{Caveats about NIR abundance-response calibration}

The model values are chemical-evolution abundance proxies, not full spectral-index inferences. The E-MILES tests use baseFe SSP spectra for flux weighting and baseline index measurement, but not abundance-variable response functions. A publication-level comparison to Eftekhari-style $[\mathrm{Si}/\mathrm{Fe}]$, \mgfe, and \simg\ needs response-calibrated SiI and MgI NIR modeling at the appropriate velocity dispersion and aperture.
